\chapter*{怎样阅读本册}

本册默认目标平台是本地 Linux 机器。读者可以使用普通 x86-64 Linux 台式机、笔记本或服务器完成完整实验；Termux 里可用的 Linux-like 环境适合阅读、构建小模型和运行基础测试，但不作为本册性能结论的标准平台。涉及 perf、NUMA、线程亲和性和 CPU 计数器时，不同环境能力会有差异，实验报告必须写清楚限制。

阅读时不要跳过前半本。AI 算子开发的难点不只是公式，而是公式如何落在真实硬件和软件运行时上。缓存 miss、TLB miss、false sharing、锁竞争、内存带宽饱和、线程切分过碎、数据布局不匹配，都会让看似正确的算子跑不快。先读现代 CPU 和并发部分，是为了让后面的 AI 内容有工程根基。

本册的贯穿项目叫 Linux CPU Quantized Inference。每读完一部分，就给项目增加一个可运行模块。第一部分完成工程骨架、张量容器和 benchmark 框架；第二部分加入线程池、同步原语和 SIMD 友好的循环；第三部分补 AI 基础，理解模型、张量、权重、激活、推理图和常见算子；第四部分实现 GEMM、归一化、Softmax 等 CPU 内核；第五部分加入 int8 量化权重、模型文件、推理执行器和端到端测试。

\begin{mentalmodel}
这本书的阅读顺序就是项目的成长顺序：先知道机器，后组织并发，再理解 AI 模型，然后写算子，最后把算子接成推理引擎。
\end{mentalmodel}

每章阅读时建议做三件事。第一，画出本章机制的边界，例如 cache line 管什么、锁保护什么、张量 stride 表达什么。第二，把机制放进贯穿项目，问它会影响哪个模块。第三，写一个最小实验验证它，例如改变线程数、改变矩阵布局、改变量化 scale 或改变 batch size。只读文字不跑实验，很难真正掌握性能工程。

\begin{warningbox}
不要把 AI 部分理解成“调用现成框架”。本册要训练的是算子开发和推理引擎实现能力。即使用小模型、小算子和教学格式，也必须保持真实工程的基本约束：输入可验证、模型可加载、结果可检查、性能可测量。
\end{warningbox}
